\appendix

%%%%%%%%%%%%%%%%%%%%%%%%%%%%%%%%%%%%%%%%%%%%%%%%%%%%%%%%%%%%%%
\part*{Anhang}
\addcontentsline{toc}{part}{Anhang}

\section*{Einleitung}

Der Anhang ergänzt die Hauptbände der Projective Structure Theory (PST)
durch mathematische Definitionen, vollständige Tabellen, historische
Entwicklungsschritte und Referenzmaterial.

Während die Haupttexte den argumentativen Aufbau der Theorie darstellen,
dient der Anhang ausschließlich als Nachschlagewerk.

Er verfolgt vier Ziele:

\begin{enumerate}

\item
Vollständige Dokumentation aller mathematischen Definitionen.

\item
Reproduzierbarkeit sämtlicher numerischer Untersuchungen.

\item
Historische Nachvollziehbarkeit der Entwicklung der Theorie.

\item
Einheitliche Referenz aller verwendeten Begriffe.

\end{enumerate}

Der Anhang wird kontinuierlich erweitert und bildet den technischen
Unterbau der gesamten PST.

\vspace{1cm}

%%%%%%%%%%%%%%%%%%%%%%%%%%%%%%%%%%%%%%%%%%%%%%%%%%%%%%%%%%%%%%

\tableofcontents

\newpage

%%%%%%%%%%%%%%%%%%%%%%%%%%%%%%%%%%%%%%%%%%%%%%%%%%%%%%%%%%%%%%
\part{A. Vollständige Formelsammlung}

\chapter{A.1 Ontophänomenalismus}

Alle Axiome

Alle Definitionen

Alle logischen Ableitungen

Alle verwendeten Symbole

%%%%%%%%%%%%%%%%%%%%%%%%%%%%%%%%%%%%%%%%%%%%%%%%%%%%%%%%%%%%%%

\chapter{A.2 Projective Structure Theory}

Vollständige mathematische Definitionen

Projektionsoperatoren

Informationsfunktionale

Gram-Matrizen

Kernel

Diffusionsoperatoren

Eigenwertdefinitionen

Stabilitätsfunktional

Effektive Dimension

Spektrallücke

Lorentz-Kriterien

%%%%%%%%%%%%%%%%%%%%%%%%%%%%%%%%%%%%%%%%%%%%%%%%%%%%%%%%%%%%%%
\part{B. Definition aller verwendeten Größen}

\chapter{Symbolverzeichnis}

Alphabetische Definition sämtlicher Größen

\begin{tabular}{ll}

$\mathcal P$
&
Projektionsoperator
\\

$R$
&
Mutual-Information-Kernel
\\

$d_{\mathrm{eff}}$
&
Effektive Dimension
\\

$\mathcal S$
&
Stabilitätsfunktional
\\

$\lambda_i$
&
Eigenwerte
\\

$\mu_i$
&
Spektrale Eigenwerte
\\

$\Gamma$
&
Projektionsgraph
\\

...

\end{tabular}

%%%%%%%%%%%%%%%%%%%%%%%%%%%%%%%%%%%%%%%%%%%%%%%%%%%%%%%%%%%%%%
\part{C. Tabellen sämtlicher numerischer Ergebnisse}

\chapter{Simulationen}

Für jede Simulation werden dokumentiert:

\begin{itemize}

\item
Testnummer

\item
Datum

\item
Ziel

\item
Parameter

\item
Code-Version

\item
Numerische Resultate

\item
Interpretation

\item
Konsequenzen

\end{itemize}

Geplant:

\begin{center}

TEST 001

...

TEST 002

...

...

TEST 037

...

\end{center}

%%%%%%%%%%%%%%%%%%%%%%%%%%%%%%%%%%%%%%%%%%%%%%%%%%%%%%%%%%%%%%
\part{D. Glossar}

Alphabetisches Nachschlagewerk sämtlicher PST-Begriffe.

Beispiele:

\begin{description}

\item[Attraktor]

...

\item[Bestimmtheit]

...

\item[Differenz]

...

\item[Emergenz]

...

\item[Existenz]

...

\item[OP]

...

\item[OPP]

...

\item[Projektionsphase]

...

\item[Relativität der Existenz]

...

\item[Selbstreferenz]

...

\end{description}

%%%%%%%%%%%%%%%%%%%%%%%%%%%%%%%%%%%%%%%%%%%%%%%%%%%%%%%%%%%%%%
\part{E. Literatur}

Unterteilt nach:

\begin{enumerate}

\item
Philosophie

\item
Mathematik

\item
Physik

\item
Informationstheorie

\item
Relationale Ontologien

\item
Emergenz

\item
Komplexe Systeme

\item
Projektive Geometrie

\item
Diffusion Maps

\item
Kernel PCA

\end{enumerate}

%%%%%%%%%%%%%%%%%%%%%%%%%%%%%%%%%%%%%%%%%%%%%%%%%%%%%%%%%%%%%%
\part{F. Chronologie der Entwicklung der PST}

Jede wesentliche Idee erhält:

\begin{tabular}{lll}

Datum

&
Version

&
Beschreibung

\end{tabular}

Beispielsweise:

\begin{itemize}

\item
Entstehung des OP

\item
Einführung des OPP

\item
Relativität der Existenz

\item
Projektionshypothese

\item
Emergente Raumzeit

\item
Projektionsphasen

\item
Fake-5D-Hypothese

...

\end{itemize}

%%%%%%%%%%%%%%%%%%%%%%%%%%%%%%%%%%%%%%%%%%%%%%%%%%%%%%%%%%%%%%
\part{G. Chronologie aller numerischen Tests}

Hier werden sämtliche Simulationen in zeitlicher Reihenfolge dokumentiert.

Jeder Eintrag enthält:

\begin{enumerate}

\item
Testnummer

\item
Motivation

\item
Fragestellung

\item
Skript

\item
Parameter

\item
Ergebnisse

\item
Interpretation

\item
Auswirkungen auf die Theorie

\end{enumerate}

%%%%%%%%%%%%%%%%%%%%%%%%%%%%%%%%%%%%%%%%%%%%%%%%%%%%%%%%%%%%%%
\part{H. Offene mathematische Vermutungen}

Hier werden ausschließlich ungelöste Probleme gesammelt.

Beispiele:

\begin{enumerate}

\item
Existiert ein universeller Projektionsoperator
$\mathcal P^\ast$?

\item
Warum entsteht die Lorentz-Signatur?

\item
Lässt sich der Projektionsoperator aus den OP-Axiomen
herleiten?

\item
Warum besitzt Diffusion Maps eine besondere Stabilität?

\item
Existiert ein Variationsprinzip für Projektionen?

\item
Wie entstehen Eichsymmetrien aus Projektionsprozessen?

\item
Kann Gravitation vollständig als Projektionsphänomen
verstanden werden?

\item
Existieren mehrere stabile Projektionsphasen?

\item
Ist die beobachtete Quantisierung ausschließlich
projektionsbedingt?

\item
Wie entsteht Zeit aus projektiver Dynamik?

\end{enumerate}

%%%%%%%%%%%%%%%%%%%%%%%%%%%%%%%%%%%%%%%%%%%%%%%%%%%%%%%%%%%%%%

\section*{Schlussbemerkung}

Der Anhang bildet das technische Gedächtnis der Projective Structure Theory.

Während sich die Hauptbände im Laufe der Forschung weiterentwickeln,
bleibt der Anhang das vollständige Archiv aller mathematischen,
numerischen und historischen Grundlagen.

Er gewährleistet damit die vollständige Nachvollziehbarkeit der
Entwicklung der Theorie und ermöglicht die Reproduktion sämtlicher
Ergebnisse.

